% VAL-011 -- generado por herramientas/gen_casos.py; editar el texto libremente.
\section{VAL-011: Voladizo espacial: flexión biaxial y torsión de Saint-Venant}\label{sec:VAL-011}
\begin{ficha}
Fuente & [S] ($GJ/L$, $PL^3/3EI$)\\
Ticket & MAM-100\\
Engine & commit \texttt{\EngineCommit}\\
Binario & SHA-256 \texttt{\EngineBinSHA}\\
Entrada & SHA-256 \texttt{\EntradaSHA{VAL-011}}\\
Tolerancia & $|c-r|\le A+R|r|$, $A=0$, $R=1e-08$ (clase consolidada; orden de ejecución DESCONOCIDO; el marco también distingue preregistro, DESCONOCIDO y VAL-037 tras la r1 de \#23)\\
Estado de datos & \directlua{mamut.estado('VAL-011')}\\
\end{ficha}

\paragraph{Modelo.}
\begin{center}\diag{voladizo}\end{center}

\paragraph{Referencia y derivación.} Derivación propia pendiente de redactar a partir del fixture del caso (provenance del Engine). % DERIV-PENDIENTE


\paragraph{Reproducción.}
\reproducir{VAL-011}

% Q2-PLACEHOLDER INICIO (VAL-011)
% Solo Q2 escribe datos/VAL-011/opensees.csv y el texto de modelo OpenSees de este bloque,
% en su propio PR. J2 no lee ni escribe valores de OpenSees.
\paragraph{Modelo de contraste (OpenSees).} \QdosPendiente
% Q2-PLACEHOLDER FIN (VAL-011)

\paragraph{Resultados.}
\begin{center}
\begin{tikzpicture}
\begin{axis}[grafico, xlabel={$x$}, ylabel={$y$}]
\directlua{mamut.serie('VAL-011','referencia','x','y','ref','Referencia analítica')}
\directlua{mamut.serie('VAL-011','engine','x','y','eng','MAMUT Engine')}
\directlua{mamut.serie('VAL-011','opensees','x','y','ops','OpenSees (contraste)')}
\end{axis}
\end{tikzpicture}
\end{center}
\begin{center}\small
\begin{tabular}{rrrrrrl}\toprule
$x$ & Ref. & Engine & OpenSees & Error rel. & Tol. & Veredicto\\\midrule
\directlua{mamut.tabla('VAL-011','x','y',0,1e-08)}
\bottomrule\end{tabular}
\end{center}
